%% ma: L12-B18-0008
%% lop: 12
%% chuong: 6
%% bai: 18
%% dang: 12.18.3
%% loai: TLN
%% muc: VD
%% dap_an: "0,6"
%% nguon: "(NBV)-ĐỀ SỐ 17-MỖI NGÀY 1 ĐỀ THI 2025.pdf (D17, MỖI NGÀY 1 ĐỀ - Đề số 17), Phần III Câu 3"
%% kien_thuc: "Bài 18 (xác suất có điều kiện), Lớp 11 Bài 29 (công thức cộng xác suất)"
%% trang_thai: chinh-thuc
%% ghi_chu: "Đáp án gốc 0,6 đúng (P(II ∩ không I) = 0,76 - 0,4 = 0,36, P(không I) = 0,6). Trùng ý tưởng với dạng đã duyệt 12.18.3 (cho P(A), P(A∪B), tính xác suất có điều kiện), nhập chính thức. Dữ kiện 0,26 là thừa. Bỏ ảnh trang trí."
\begin{ex}
Một công ty tiến hành hai dự án I và II. Biết rằng xác suất thành công của ít nhất một trong hai dự án là $0{,}76$; xác suất thành công của dự án I là $0{,}4$; xác suất để dự án I thành công nhưng dự án II không thành công là $0{,}26$. Tìm xác suất để dự án II thành công nếu biết rằng dự án I không thành công.
\shortans{0,6}
\loigiai{Gọi $A$, $B$ lần lượt là biến cố dự án I, dự án II thành công. Ta có $P(A\cup B)=0{,}76$, $P(A)=0{,}4$, $P(\overline A)=0{,}6$.

$A\cup B$ là hợp của hai biến cố xung khắc $A$ và $\overline A\cap B$, nên $P(\overline A\cap B)=P(A\cup B)-P(A)=0{,}76-0{,}4=0{,}36$.

Vậy $P(B\mid\overline A)=\dfrac{P(\overline A\cap B)}{P(\overline A)}=\dfrac{0{,}36}{0{,}6}=0{,}6$.}
\end{ex}
